% A uniform finite membership reduction; no full descent lift is inferred.
\subsection{A fixed finite test for the six simultaneous source words}

Retain the literal $\varphi_P:F\twoheadrightarrow B$ of
(B.Q5).  For $b\in\{1,s,r,r^2,sr,sr^2\}$, set
\[
 \varphi_b=\varphi_P\circ\operatorname{Ad}_b|_F,
 \qquad N_b=\ker\varphi_b,
 \qquad \eta=(\varphi_1,\varphi_s,\varphi_r,
                    \varphi_{r^2},\varphi_{sr},\varphi_{sr^2})
          :F\longrightarrow B^6.                              \tag{B.R1}
\]
Every coordinate is the displayed actual homomorphism from
the same source $F$.  Put $J=\eta(F)$ and define the following
normal subgroups of the target $B$:
\[
 \begin{aligned}
 T_2&=\langle\!\langle\tau^2\rangle\!\rangle_B,&
 T_3&=\langle\!\langle\tau^3\rangle\!\rangle_B,\\
 L_0&=[T_2,T_3],& L_{i+1}&=[L_i,B]\quad(0\leq i\leq3).
 \end{aligned}                                                   \tag{B.R2}
\]
All commutator and normal subgroups are closed.  Thus $L_3$
is the iterated commutator of $T_2,T_3$ and three further
copies of $B$, while $L_4=[L_3,B]$.

\begin{result}{B.2.10}{Proved}{Uniform finite tests for simultaneous word membership}
Each quotient $B/L_i$ in (B.R2) is finite, and
\[
                       L_3^6\subseteq J,
                       \qquad L_4^6\subseteq J'.                \tag{B.R3}
\]
Let $\pi_i:B^6\to(B/L_i)^6$ be the coordinatewise quotient
and $J_i=\pi_i(J)$.  Then the exact two word images are
\[
 \eta(F)=\pi_3^{-1}(J_3),\qquad
 \eta(F')=\pi_4^{-1}(J_4').                                    \tag{B.R4}
\]
In particular both membership questions are determined by
fixed finite quotients, independently of a later target
level $Q$.  They are membership tests for the six prescribed
word evaluations; (B.R4) asserts no positive automorphism
or canonical central-cover lifting conclusion.
\end{result}

\begin{proof}
The orders of the three peripheral images determine the
six kernels exactly.  Here an entry $\infty$ means that
no power relation on that peripheral is imposed.  The final
column gives the image of the actual kernel under the first
coordinate $\varphi_1=\varphi_P$:
\[
\begin{array}{c|ccc|l|c}
 b&x&y&z&N_b\text{ as a normal subgroup of }F&\varphi_1(N_b)\\ \hline
 1&2&\infty&3&\langle\!\langle x^2,z^3\rangle\!\rangle&1\\
 s&\infty&2&3&\langle\!\langle y^2,z^3\rangle\!\rangle&T_2\\
 r&\infty&3&2&\langle\!\langle y^3,z^2\rangle\!\rangle&B\\
 r^2&3&2&\infty&\langle\!\langle x^3,y^2\rangle\!\rangle&B\\
 sr&2&3&\infty&\langle\!\langle x^2,y^3\rangle\!\rangle&T_3\\
 sr^2&3&\infty&2&\langle\!\langle x^3,z^2\rangle\!\rangle&B
\end{array}                                                     \tag{B.R5}
\]
For example $\varphi_1(x)=s$, $\varphi_1(y)=\tau$ and
$\varphi_1(z)=r^{-1}$.  In the $r$ row, the two relations
give $r=1$ and $\tau^3=1$, and hence $s=1$ as well.  In the
$r^2$ row they give $s=1$ and $\tau^2=1$, and hence $r=1$.
The last row kills both $s$ and $r$.  The other two nontrivial
rows give exactly $T_2$ and $T_3$.  The table of kernels
follows from the actual conjugation formulas
$r(x,y,z)r^{-1}=(y,z,x)$ and the interchange of the $x,y$
peripheral conjugacy classes by $s$; conjugating a peripheral
does not change its closed normal power subgroup.

The quotients of $B$ by $T_2$ and $T_3$ are respectively
the triangle groups of signatures $(2,3,2)$ and $(2,3,3)$,
namely $S_3$ and $A_4$.  These identifications also follow
directly from the standard presentations: use a transposition
and a three-cycle in $S_3$, and a double transposition and
a three-cycle in $A_4$.  Their orders are six and twelve,
as is verified by the elementary coset enumerations for
these spherical triangle presentations.  The actual joint
map is onto $S_3\times A_4$.  To see this without assuming
independence, a proper subdirect product would give a common
nontrivial quotient.  The only proper nontrivial quotient
of $S_3$ is $C_2$, and that of $A_4$ is $C_3$; neither
group is a quotient of the other.  There is no such common
quotient.  Thus, with $K_0=T_2\cap T_3$,
\[
                 B/K_0\cong S_3\times A_4.                     \tag{B.R6}
\]
Equally, $T_2T_3=B$ follows immediately by imposing both
$\tau^2=1$ and $\tau^3=1$: then $\tau=1$, and the
orders of $s=r^{-1}$ force $s=r=1$.

The subgroup $K_0/[T_2,T_3]$ is central in
$B/[T_2,T_3]$.  Indeed an element of $K_0$ belongs to
both $T_2$ and $T_3$; its commutator with either factor
lies in $[T_2,T_3]$, and the two factors generate $B$.
The map from $B^{\mathrm{ab}}=C_6$ to
$(S_3\times A_4)^{\mathrm{ab}}=C_2\times C_3$ is an
isomorphism on the specified $s,r$.  Hence $K_0\subseteq B'$.
Transfer, exactly as in (B.Q1), shows that the central group
$K_0/[T_2,T_3]$ is finite: it is topologically finitely
generated, lies in the derived subgroup of its ambient
group, and is annihilated by the finite index $72$.
This proves that $L_0$ is open.  It is contained in $B'$.
For each successive $L_i$, the group
$L_i/[L_i,B]$ is central in $B/[L_i,B]$, contained in its
derived subgroup, and of finite index $|B/L_i|$.  The same
transfer argument proves it finite.  Induction proves
finiteness of every $B/L_i$ in (B.R2).

Fix one coordinate $b$ of (B.R1).  Conjugation by $b$
permutes the six kernels and carries $N_b$ to $N_1$.
Therefore the five images $\varphi_b(N_{b'})$, for
$b'\ne b$, are precisely $T_2,T_3,B,B,B$, in a suitable
order, by (B.R5).  Form the iterated commutator of these
five source normal subgroups in that order.  It lies in
their intersection, and its image under $\varphi_b$ is
$L_3$.  Its other five coordinates under $\eta$ are
trivial.  Thus $J$ contains a copy of $L_3$ supported in
each individual coordinate.  Multiplying those six
subgroups proves the first containment of (B.R3).

The projection of $J$ to each factor $B$ is onto.  Commuting
the copy of $L_3$ supported in that factor with all of $J$
therefore gives the copy of $[L_3,B]=L_4$ supported in the
same factor.  These commutators belong to $J'$, so the
second containment of (B.R3) follows.  In particular both
$J$ and $J'$ are open in $B^6$.

The first equality in (B.R4) follows because
$\ker\pi_3=L_3^6\subseteq J$.  For the second, a continuous
epimorphism maps the closed derived subgroup onto the
closed derived subgroup of its image, so $\eta(F')=J'$.
Also $\pi_4(J')=J_4'$ and
$\ker\pi_4=L_4^6\subseteq J'$.  These two facts give
exactly the second equality in (B.R4).
\end{proof}

\paragraph{Preservation of the finite test by positive target actions.}
Every positive target automorphism taking $\tau$ to a
conjugate of a unit power preserves $T_2,T_3$, and hence
each $L_i$ in (B.R2).  Thus its action on these fixed finite
quotients is an actual induced action.  Substitution of the
full frame words (B.Q11), retaining both free exponents,
gives an exact finite test for their simultaneous realization
by one commutator word.  The additional conditions needed
to produce a positive automorphism of $E_Q^{\mathrm c}$
are still the conditions in (B.Q2).  The finite-index
argument above does not remove them.
